% ============================================
% Response to Reviewer 2
% ============================================
\section*{Response to Reviewer 2}

We thank the reviewer for identifying the remaining conceptual ambiguity and for suggesting a clear way to resolve it.

\subsection*{R2-1. Distinguish the Dominance signature from assembly-history evidence}

\reviewercomment{The operational definition of community-level selection as parental-origin-correlated persistence makes the statement that Dominance ``provides outcome-level evidence for community-level selection'' potentially circular. Dominance should instead be presented as the compositional signature used to quantify origin-correlated persistence. The excess Dominance observed after separate parental-community assembly relative to direct assembly provides the stronger evidence that shared assembly history generates this origin dependence. The reviewer asks that this distinction be made explicit in the main text.}

We agree and have revised the manuscript accordingly. In \S2.1, we clarified
the distinction between Dominance and community-level selection, identifying
Dominance as the compositional signature used to quantify origin-correlated
persistence. At the end of \S2.2, we also added a substantive discussion of
assembly history and its effect on the fraction of Dominance outcomes, based on
the comparison between coalescence and direct assembly.
